\documentclass[11pt,a4paper]{article}

\usepackage[margin=25mm,headheight=14pt]{geometry}
\usepackage[T1]{fontenc}
\usepackage{lmodern}
\usepackage{microtype}
\usepackage{mathtools,amssymb,amsthm}
\usepackage{booktabs,longtable,tabularx,array}
\usepackage[table]{xcolor}
\usepackage{enumitem}
\usepackage{seqsplit}
\usepackage{fancyhdr}
\usepackage{hyperref}
\usepackage{bookmark}

\hypersetup{
  colorlinks=true,
  linkcolor=blue!45!black,
  urlcolor=blue!55!black,
  citecolor=blue!45!black,
  pdftitle={Strengthened results in the Cauchy-sequence completion of the reals},
  pdfauthor={Source audit of cubical-master-reals}
}

\definecolor{deepblue}{HTML}{17365D}
\definecolor{softblue}{HTML}{EAF2F8}
\definecolor{softgreen}{HTML}{EAF6EC}
\definecolor{softamber}{HTML}{FFF4D6}
\definecolor{softred}{HTML}{FBE9E7}
\definecolor{softgray}{HTML}{F4F5F6}

\pagestyle{fancy}
\fancyhf{}
\lhead{Cauchy-real comparison}
\rhead{Source audit, 12 August 2026}
\cfoot{\thepage}

\setlist[itemize]{leftmargin=1.5em,itemsep=.25em,topsep=.3em}
\setlist[enumerate]{leftmargin=1.7em,itemsep=.35em,topsep=.3em}
\renewcommand{\arraystretch}{1.18}
\setlength{\LTpre}{.5em}
\setlength{\LTpost}{.5em}
\setlength{\emergencystretch}{2em}

% HoTT Book-style conventions, kept self-contained for this report.
\newcommand{\N}{\mathbb{N}}
\newcommand{\Q}{\mathbb{Q}}
\newcommand{\Qp}{\mathbb{Q}_{+}}
\newcommand{\R}{\mathbb{R}}
\newcommand{\RC}{\mathbb{R}_{\mathrm C}}
\newcommand{\RP}{\mathbb{R}_{\mathrm P}}
\newcommand{\RS}{\mathbb{R}_{\mathrm S}}
\newcommand{\Type}{\mathcal{U}}
\newcommand{\Prop}{\mathsf{Prop}}
\newcommand{\brck}[1]{\left\lVert #1\right\rVert}
\newcommand{\prd}[2]{\prod_{#1} #2}
\newcommand{\sm}[2]{\sum_{#1} #2}
\newcommand{\eqv}{\simeq}
\newcommand{\close}[1]{\mathrel{\sim_{#1}}}
\newcommand{\Tail}{\mathsf{Tail}}
\newcommand{\Cch}{\mathsf{C}_{\mathrm{ch}}}
\newcommand{\Cbd}{\mathsf{C}_{\mathrm{bd}}}
\newcommand{\Crad}{\mathsf{C}_{\mathrm{rad}}}
\newcommand{\Lim}{\mathsf{Lim}}
\newcommand{\Sum}{\mathsf{Sum}}
\newcommand{\AbsC}{\mathsf{AbsC}}
\newcommand{\UC}{\mathsf{UC}}
\newcommand{\RU}{\mathsf{RU}}
\newcommand{\PS}{\mathsf{PS}}
\newcommand{\RatContr}{\mathsf{RatContr}}
\newcommand{\CP}{\mathbin{*_{\mathrm C}}}
\newcommand{\code}[1]{{\ttfamily\seqsplit{#1}}}
\newcommand{\strength}{\textcolor{deepblue}{\textsf{strengthening}}}
\newcommand{\newcov}{\textcolor{green!40!black}{\textsf{new coverage}}}
\newcommand{\gap}{\textcolor{red!55!black}{\textsf{older side broader}}}

\newtheorem{finding}{Finding}
\newtheorem{caveat}{Scope caveat}

\title{\textbf{Strengthened Results in the Cauchy-Sequence Completion of the Reals}\\
\large A theorem-level comparison with \code{HITs/CauchyReals} and \code{HITs/CauchyRealsP}}
\author{Source audit of \code{cubical-master-reals}}
\date{12 August 2026}

\begin{document}
\maketitle

\begin{abstract}
The material improvement in the \code{Completion/CauchySequence/Reals} development is not
that it proves ordinary completeness once again.  It changes the logical strength of the
input accepted by completeness and then consistently propagates that change through
sequences, series, uniform convergence, power series, and integration.  A Cauchy tail is
chosen only at the radius currently being used, under propositional truncation.  No global
modulus is manufactured.  Consequently, the new development proves the choice-free
equivalence
\[
  \Crad(x) \eqv \Lim(x),
\]
whereas the older sequential APIs start from a simultaneously chosen modulus.  This yields
genuinely stronger theorem interfaces: arbitrary-real rather than rational scaling of
Cauchy series, sharp hypotheses such as $|x|<1$ rather than a selected rational contraction,
Mertens theorems without selected sum bounds, and uniform-limit/integration results without
global modulus data.  The report also records an important negative result: a suitably
structured inverse from the regular completion to the new completion is equivalent to the
missing sequential tail-choice principle.  Finally, the older \code{HITs/CauchyReals}
development still has substantially broader global calculus, logarithm, $\pi$, inverse
trigonometry, circle, and complex-number coverage.  Thus the new line is logically sharper
in its sequential interfaces, not yet globally dominant in library coverage.
\end{abstract}

\tableofcontents

\section{Scope, names, and the comparison being made}

There are three related, but not interchangeable, code lines in scope.

\paragraph{$\RC$ --- the HoTT-style higher inductive-inductive real.}
  \path{Cubical/HITs/CauchyReals}.  Its constructors follow HoTT Book
  Definition~11.3.2: rational points, limits of $\Qp$-indexed Cauchy
  approximations, and paths identifying points close at every radius.

\paragraph{$\RP$ --- the generic regular completion specialised to $\Q$.}
  \path{Cubical/HITs/CauchyRealsP}.  This is the direct predecessor of most of
  the sequence/series/power-series code audited here.  Its foundation re-exports
  \path{Relation/Premetric/Completion/Base}; its analysis API ordinarily exposes
  chosen sequential moduli.

\paragraph{$\RS$ --- the radiuswise Cauchy-sequence completion.}
  \path{Cubical/Relation/Premetric/Completion/CauchySequence/Reals}.  Its limit
  constructor takes an $\N$-sequence and a separately truncated tail at each
  radius.  Most direct old/new theorem pairs are therefore $\RP$ versus $\RS$.

The source contains an unconditional canonical nonexpanding map
\[
  c : \RP \longrightarrow \RS
\]
which fixes rational points; companion modules prove preservation of distance, lattice,
multiplication, and commutative-ring structure.  No unconditional inverse is asserted.
There is also no source theorem identifying the full \code{HITs/CauchyReals} development
with $\RS$.  Accordingly, this report does not silently transport old theorems across an
unproved equivalence.

We use HoTT Book conventions: $\Qp$ for positive rationals, $\brck{A}$ for propositional
truncation, dependent products and sums for $\prod$ and $\sum$, and $\eqv$ for equivalence.
Write $x_m\close{\varepsilon}x_n$ for premetric closeness.  This notation suppresses
the implementation's two concrete closeness relations, which is immaterial to the logical
comparison.

\paragraph{Classification used below.}
A \strength{} weakens a premise, strengthens a conclusion, or removes data that cannot be
recovered constructively.  \newcov{} means a construction is new relative to $\RP$ but is
not necessarily a stronger counterpart of a theorem already in $\RC$.  A \gap{} marks an
area where the mature $\RC$ hierarchy remains broader.

\section{The foundational material difference}

\subsection{The HoTT constructor and the new constructor}

The HoTT Book and \code{HITs/CauchyReals} use Cauchy approximations.  Schematically,
\begin{equation}
\label{eq:hott-lim}
  \mathsf{lim}_{\mathrm C} :
  \prd{x:\Qp\to\RC}
  \left(\prd{\delta,\varepsilon:\Qp}
    x_\delta\close{\delta+\varepsilon}x_\varepsilon\right)
  \longrightarrow \RC .
\end{equation}
This is the free-completion idea of HoTT Book \S11.3: limits and quotient paths are added
simultaneously, avoiding a choice of representatives for a sequence in a quotient.

The new completion instead has a genuinely sequential constructor
\begin{equation}
\label{eq:seq-lim}
  \mathsf{lim}_{\mathrm S} :
  \prd{x:\N\to\RS} \Crad(x) \longrightarrow \RS,
\end{equation}
where
\begin{align}
  \Tail(x,\varepsilon)
  &\;\coloneqq\;
  \sm{N:\N}
  \prd{m,n:\N}
  (N\le m)\to(N\le n)\to(x_m\close{\varepsilon}x_n),
  \label{eq:tail}\\
  \Crad(x)
  &\;\coloneqq\;
  \prd{\varepsilon:\Qp}\brck{\Tail(x,\varepsilon)}.
  \label{eq:radiuswise}
\end{align}
The one common cutoff needed at a particular $\varepsilon$ exists merely; the definition
contains no function selecting all cutoffs at once.

\subsection{Three placements of truncation}

The comparison module isolates three distinct propositions/data types:
\begin{align}
  \Cch(x)
  &\;\coloneqq\;
    \prd{\varepsilon:\Qp}\Tail(x,\varepsilon)
    \;\eqv\;
    \sm{M:\Qp\to\N}
    \prd{\varepsilon:\Qp}\text{tail-law}(M\varepsilon),
    \label{eq:chosen}\\
  \Cbd(x)
  &\;\coloneqq\; \brck{\Cch(x)},
    \label{eq:bundled}\\
  \Crad(x)
  &\;\coloneqq\;
    \prd{\varepsilon:\Qp}\brck{\Tail(x,\varepsilon)}.
    \label{eq:radagain}
\end{align}
There are constructive maps
\[
  \Cch(x)\longrightarrow\Cbd(x)\longrightarrow\Crad(x),
\]
but the missing reverse direction is exactly the passage
\begin{equation}
\label{eq:choice-gap}
  \left(\prd{\varepsilon:\Qp}\brck{A_\varepsilon}\right)
  \longrightarrow
  \brck{\prd{\varepsilon:\Qp}A_\varepsilon}.
\end{equation}
That is a countable-choice principle over $\Qp$, not a law of propositional truncation.
The older $\RC$ sequential predicate \code{IsCauchySequence'} and $\RP$ predicate
\code{isCauchySeq} inhabit the chosen side of this divide (their concrete packaging differs,
but both expose simultaneous tail choices).  The $\RS$ predicate inhabits the radiuswise side.

\begin{finding}[Exact obstruction]
Let $\mathsf{TC}$ be the real-specific sequential tail-choice statement that turns every
radiuswise Cauchy sequence in the intermediate sequential presentation into one bundled
sequential witness.  The source proves
\begin{equation}
\label{eq:structured-inverse}
  \mathsf{StructuredInverse}(c) \eqv \mathsf{TC}.
\end{equation}
Here ``structured'' means an equivalence $c:\RP\eqv\RS$ whose inverse also supplies the
bundled Cauchy witness demanded by the regular/sequential presentation.  Propositional choice
over $\Qp$ implies $\mathsf{TC}$ and hence such an inverse.  Bare invertibility of $c$ is
deliberately not claimed sufficient, because inverse values alone do not compute the missing
tail witnesses.
\end{finding}

Equation~\eqref{eq:structured-inverse} is the sharpest material comparison in the repository.
It says where choice would enter and characterises exactly the extra structure it would buy.
This directly echoes the warning in HoTT Book \S11.3: quotient-first Cauchy completions need
countable choice to lift a sequence of quotient elements to representatives.  The new issue
is not representative choice but tail choice; the same placement-of-truncation phenomenon is
made explicit in the type.

\section{The principal strengthened theorem}

For a sequence $x:\N\to\RS$, define its type of limits by
\begin{align}
  \mathsf{isLim}(x,\ell)
  &\;\coloneqq\;
  \prd{\varepsilon:\Qp}
  \brck{\sm{N:\N}\prd{n:\N}
    (N\le n)\to(x_n\close{\varepsilon}\ell)},\\
  \Lim(x)
  &\;\coloneqq\;
  \sm{\ell:\RS}\mathsf{isLim}(x,\ell).
\end{align}
Separatedness makes $\Lim(x)$ a proposition.  The constructor gives
$\Crad(x)\to\Lim(x)$; the triangle inequality gives
$\Lim(x)\to\Crad(x)$ without choosing tails at all radii.  Therefore:
\begin{equation}
\boxed{\qquad \Crad(x)\;\eqv\;\Lim(x). \qquad}
\label{eq:main-equivalence}
\end{equation}
In source this is \code{isCauchySequence}$\simeq$\code{sequenceLimit}.  The older direct
counterpart \code{cauchySequenceLimit} proves only
\[
  \Cch(x)\longrightarrow\Lim(x).
\]
Thus equation~\eqref{eq:main-equivalence} is stronger in both directions: its forward map
accepts less structured Cauchy evidence, and its reverse map establishes a full criterion for
convergence.  It does \emph{not} imply $\Crad(x)\to\Cbd(x)$; completeness returns a real limit,
not a global modulus.

Two important corollaries are then obtained at the correct strength:
\begin{align*}
  &f\text{ continuous},\ \Crad(x)
    \quad\Longrightarrow\quad \Crad(f\circ x),\\
  &f\text{ nonexpanding},\ \mathsf{isLim}(x,\ell)
    \quad\Longrightarrow\quad \mathsf{isLim}(f\circ x,f\ell).
\end{align*}
The first implication uses completeness of $\RS$: obtain the unique limit, apply sequential
continuity, then use convergence $\Rightarrow$ radiuswise Cauchy.  It would be false for an
arbitrary incomplete target space, so the proof is not merely a repackaged Lipschitz lemma.

\section{Theorem-by-theorem strength ledger}

\rowcolors{2}{softgray}{white}
\begin{longtable}{>{\raggedright\arraybackslash}p{.17\textwidth}
                  >{\raggedright\arraybackslash}p{.32\textwidth}
                  >{\raggedright\arraybackslash}p{.39\textwidth}}
\toprule
\textbf{Area / class} & \textbf{Earlier statement or burden} &
\textbf{Strengthened $\RS$ result and why it matters}\\
\midrule
\endfirsthead
\toprule
\textbf{Area / class} & \textbf{Earlier statement or burden} &
\textbf{Strengthened $\RS$ result and why it matters}\\
\midrule
\endhead
\bottomrule
\endfoot

Cauchy predicate
\newline\strength
& $\Cch(x)$: a cutoff at every radius is simultaneously selected.
& $\Crad(x)=\prod_\varepsilon\brck{\Tail(x,\varepsilon)}$.  This is implied by the old
predicate; the converse is precisely tail choice and is not assumed.\\

Completeness
\newline\strength
& \code{cauchySequenceLimit}: $\Cch(x)\to\Lim(x)$.
& \code{isCauchySequence}$\simeq$\code{sequenceLimit}:
$\Crad(x)\eqv\Lim(x)$.  This is the central stronger result, not merely a new constructor.\\

Eventually equal and ordered
\newline\strength
& The established predicates store $\sum_N\prod_{n\ge N}(\cdots)$ and expose $N$.
& \code{EventuallyEqualSequence}, \code{EventuallyLeSequence}, and
\code{EventuallyLtSequence} truncate the cutoff.  Eventual $\le$ and even eventual $<$ pass
to $\le$ of limits without choosing a witness.\\

Asymptotic closeness
\newline\strength
& A selected cutoff is returned at each radius and clients often retain the witness.
& \code{EventuallyClose} follows the same radiuswise truncation discipline and preserves
Cauchyness and limits by proposition-valued elimination.\\

Series criterion
\newline\strength
& $\mathsf{CauchySeries}(a)\to\mathsf{Summable}(a)$, with the former built from a chosen
modulus of partial sums.
& \code{isCauchySeries}$\simeq$\code{summable} gives an equivalence at the radiuswise
strength:
\[
  \Crad\!\left(n\mapsto\sum_{k<n}a_k\right)\eqv\Sum(a).
\]
No global series modulus is recovered or required.\\

Scalar closure
\newline\strength
& \code{rationalScaleSeriesIsCauchy} and
\code{rationalScaleSeriesIsAbsolutelyCauchy} only scale by $q:\Q$.
& \code{scaleSeriesIsCauchy} and \code{scaleSeriesIsAbsolutelyCauchy} accept every
$c:\RS$.  Native completeness plus continuity of multiplication removes the rational-only
restriction.\\

Majorants
\newline\strength
& From $|a_n|\le b_n$, old power-series theorems additionally request
$0\le b_n$ for all $n$.
& \code{coefficientMajorantIsNonnegative} derives
\[
  0\le |a_n|\le b_n,
\]
so all coefficient-majorant convergence theorems drop the redundant nonnegativity premise.\\

Absolute bounds
\newline\strength
& Mertens and iterated-product clients select explicit $A:\Qp$ bounding an absolute sum.
& \code{mereAbsoluteCauchySeriesSumBound} proves only
\[
  \brck{\sm{A:\Qp}\mathsf{AbsSumBound}(A,a)}.
\]
Clients open such a bound only inside a proposition.  The useful conclusion survives while
the arbitrary bound is removed from the public theorem.\\

Mertens product
\newline\strength
& Old absolute Mertens statements take $A,B:\Qp$ and proofs that the two absolute sums are
bounded by them; separate theorems establish ordinary or absolute Cauchyness.
& For the Cauchy product $(a\CP b)_n=\sum_{k\le n}a_kb_{n-k}$:
\[
 \begin{gathered}
 \AbsC(a)\to\AbsC(b)\to\AbsC(a\CP b),\\[-2pt]
 \sum(a\CP b)=(\sum a)(\sum b).
 \end{gathered}
\]
Theorems \code{absoluteMertensCauchyProductIsAbsolutelyCauchy} and
\code{absoluteMertensSeriesSum-cauchyProduct} need no selected sum bounds.\\

Geometric series
\newline\strength
& Callers provide
\[
 \begin{aligned}
 \RatContr(x)\coloneqq\sum_{q:\Qp}&(|x|<q)\\[-2pt]
 &{}\times(q<1).
 \end{aligned}
\]
& \code{geometricSeriesIsAbsolutelyCauchyOfAbs<1} and its summability/value variants assume
only the sharp open condition $|x|<1$.  Density yields $\brck{\RatContr(x)}$, which is opened
only into the propositional conclusion.\\

Bounded power series
\newline\strength
& With $|a_n|\le C$, convergence at $x$ requires a selected
$\RatContr(x)$.
& \code{powerSeriesIsAbsolutelyCauchyOfAbs<1},
\code{powerSeriesIsAbsolutelySummableOfAbs<1}, and
\code{powerSeriesIsSummableOfAbs<1} require only $|x|<1$.  Product and interior convergence
layers reuse the same sharp interface rather than leaking an intermediate rational radius.\\

Uniform series
\newline\strength
& A whole uniform modulus $M:\Qp\to\N$ is part of
\code{UniformlyCauchySeries}.
& The new definition is
\[
 \begin{gathered}
 \mathsf{UCauchy}(S)\coloneqq\\[-2pt]
 \displaystyle\prod_{\varepsilon:\Qp}
 \brck{\mathsf{UniformTail}(S,\varepsilon)}.
 \end{gathered}
\]
Here \(\mathsf{UniformTail}\) contains one cutoff common to all parameters.  Thus the
definition does not collapse to pointwise convergence.  It yields uniform convergence to
the canonical sum and continuity of that sum with no chosen global modulus.\\

Uniform continuity
\newline\strength
& Conventional APIs package a modulus choosing an input radius for every requested output
radius.
& \code{IsRadiuswiseUniformlyContinuous} uses
\[
 \RU(f)\coloneqq\prod_{\varepsilon:\Qp}
 \brck{\mathsf{ModulusAt}(f,\varepsilon)}.
\]
Here \(\mathsf{ModulusAt}\) contains an adequate input radius and its preservation law.
Conventional uniform continuity implies $\RU(f)$; no converse is exported because it would
again commute product with truncation.\\

Uniform limits and integration
\newline\strength{} / \newcov
& The older $\RP$ series hierarchy has no Riemann-integration layer; the older $\RC$ layer
uses stronger uniform-continuity/modulus interfaces.
& An arbitrary uniformly convergent sequence of radiuswise uniformly continuous functions
has a radiuswise uniformly continuous limit.  The completed dyadic integral preserves both
global and interval-local uniform limits:
\[
 \begin{gathered}
 f_n\to f\text{ uniformly on }[a,b]\\[-2pt]
 \Longrightarrow\ \int_a^b f_n\to\int_a^b f.
 \end{gathered}
\]
Completing the sequence of integrals independently is proved equal to integrating the limit.\\

Transcendental construction
\newline\newcov
& $\RP$ stops at generic power-series machinery.  $\RC$ already has a different, mature
transcendental hierarchy.
& $\RS$ now constructs power-series $\exp$, $\sin$, $\cos$, hyperbolic functions, a partial
tangent layer, local arctangent/reciprocal results, and $-\log(1-x)$ on $|x|<1$, with addition,
derivative, order, and identity theorems.  This is substantial new coverage relative to
$\RP$, but not by itself a stronger theorem than the corresponding mature $\RC$ results.\\
\end{longtable}
\rowcolors{2}{}{}

\section{Four representative proof improvements}

\subsection{Sharp open hypotheses without selected rational witnesses}

The geometric-series change is the cleanest recurring proof pattern.  The old implementation
already proves convergence from $\RatContr(x)$.  Rational density proves only
\[
  |x|<1\quad\longrightarrow\quad\brck{\RatContr(x)},
\]
which is exactly what constructive density should provide.  In $\RS$, absolute Cauchyness is
a proposition, so truncation elimination gives
\[
  |x|<1\quad\longrightarrow\quad
  \AbsC\bigl(n\mapsto x^n\bigr)
\]
without selecting $q$.  The same proof architecture supplies the sharp hypotheses for bounded
power series and their products.  The mathematical condition visible to a client is now the
standard one; the rational $q$ remains an internal estimate.

This is materially stronger than adding a convenience lemma only if the target is genuinely
propositional.  The report therefore does not claim a construction of a canonical contraction
radius.  Such a construction is neither present nor needed.

\subsection{Mertens without bound plumbing}

For an absolutely Cauchy series $a$, Archimedean boundedness gives
\[
  \AbsC(a)\longrightarrow
  \brck{\sm{A:\Qp}\mathsf{AbsSumBound}(A,a)}.
\]
The old Mertens interface made $A$ and its proof caller-visible.  The new proof eliminates
mere bounds for $a$ and $b$ directly into the proposition that rectangular partial sums
approximate triangular Cauchy-product sums.  It then derives, in order,
\begin{align*}
  \AbsC(a),\AbsC(b)
    &\Longrightarrow \mathsf{RectangularApproximation}(a,b),\\
    &\Longrightarrow \mathsf{Cauchy}(a\CP b),\\
    &\Longrightarrow \AbsC(a\CP b),\\
    &\Longrightarrow \mathsf{AbsolutelySummable}(a\CP b).
\end{align*}
The public theorem has the textbook premise, while the quantitative estimates still use
explicit bounds locally.  This is a good example of truncation serving abstraction rather
than erasing information prematurely.

\subsection{Arbitrary-real scaling from native completeness}

The old Cauchy proof scaled a modulus by a rational and therefore exported only
\[
  q:\Q,\ \mathsf{CauchySeries}(a)
  \Longrightarrow \mathsf{CauchySeries}(q a).
\]
In $\RS$, a constant real sequence is radiuswise Cauchy, multiplication preserves native
Cauchy sequences, and the partial sums of $ca$ are eventually equal to
$c$ times the partial sums of $a$.  Hence
\[
  c:\RS,\ \mathsf{CauchySeries}(a)
  \Longrightarrow \mathsf{CauchySeries}(ca),
\]
with the analogous absolute result.  This is a genuine extension of the scalar domain, not
only a change in witness packaging.

\subsection{Uniformity is retained, global choice is not}

The location of truncation in the new uniform-series predicate is deliberate:
\[
  \prd{\varepsilon}\brck{\sm{N}
  \prd{x}\prd{m,n\ge N}S_m(x)\close{\varepsilon}S_n(x)}.
\]
Placing $\prod_x$ \emph{inside} the truncation retains a single cutoff common to the whole
family.  Placing it outside would yield only pointwise Cauchyness.  Placing the outer
$\prod_\varepsilon$ inside one truncation would restore a globally chosen modulus.  The new
definition occupies the exact middle position needed by uniform-limit proofs.

At a requested output radius, the continuity proof opens one common approximation index and
one input radius for that single approximant.  Both eliminations land in propositions.  This
is the same radius-by-radius discipline as the completion constructor, propagated through a
nontrivial analytic construction.

\section{Integration: what is stronger and what is only newer}

The $\RS$ integration layer contains two distinct accomplishments.

\begin{enumerate}
\item \textbf{Stronger interface discipline.}  Riemann resampling and dyadic refinement use
  $\RU(f)$, not a selected uniform modulus.  Vanishing-mesh quadrature sequences are proved
  radiuswise Cauchy and completed directly.  Integrability, uniqueness, linearity, endpoint
  agreement, and scheme independence are proposition-valued where possible.

\item \textbf{New completion-level uniform-limit construction.}  For
  $f_n,f:\RS\to\RS$, define uniform convergence by
  \[
    \prd{\varepsilon:\Qp}\brck{\sm{N:\N}
      \prd{x:\RS}\prd{n\ge N}f_n(x)\close{\varepsilon}f(x)}.
  \]
  If every $f_n$ is radiuswise uniformly continuous, the limit $f$ is too, and
  \[
    \mathsf{isLim}\!\left(
      n\mapsto I_{a,b}(f_n),\ I_{a,b}(f)\right)
  \]
  holds.  A local theorem assumes convergence only on $[a,b]$; it uses a merely existing
  rational bound on $b-a$ inside the proposition.  A second construction first completes
  the Cauchy sequence $I_{a,b}(f_n)$ and proves that independent completion agrees with
  $I_{a,b}(f)$.
\end{enumerate}

These are stronger than the $\RP$ situation, which has no integration layer, and their
hypotheses are constructively sharper than a chosen-modulus formulation.  They do not yet
subsume the mature $\RC$ calculus.  In particular, no analogue of
\path{HITs/CauchyReals/Calculus/Integration/FundamentalTheorem.agda} was found in $\RS$.

\section{Coverage comparison: do not confuse sharpness with maturity}

\begin{caveat}
The new completion should not be described as simply ``stronger than
\code{HITs/CauchyReals}.''  Its \emph{sequential propositions and theorems that consume them}
are more constructively economical.  The older $\RC$ hierarchy nevertheless contains many
major developments that have not been transported or reconstructed over $\RS$.
\end{caveat}

\rowcolors{2}{softgray}{white}
\begin{longtable}{>{\raggedright\arraybackslash}p{.20\textwidth}
                  >{\raggedright\arraybackslash}p{.34\textwidth}
                  >{\raggedright\arraybackslash}p{.34\textwidth}}
\toprule
\textbf{Topic} & \textbf{$\RC$ / $\RP$ status} & \textbf{$\RS$ status}\\
\midrule
\endfirsthead
\toprule
\textbf{Topic} & \textbf{$\RC$ / $\RP$ status} & \textbf{$\RS$ status}\\
\midrule
\endhead
\bottomrule
\endfoot

Sequential completeness
& $\RC$ and $\RP$ complete chosen-modulus sequences.
& Radiuswise Cauchy $\eqv$ convergent; exact tail-choice obstruction for the regular comparison.\\

Series and power series
& $\RP$ has a large generic hierarchy, but rational contractions, selected bounds, and
chosen uniform moduli appear in public hypotheses.
& Ported and extended with radiuswise predicates, sharp $|x|<1$ clients, arbitrary-real
scaling, bound-free Mertens, and local/uniform differentiation.\\

Differential calculus
& $\RC$ has derivative, inverse, mean-value, scaling, uniqueness, and power rules;
$\RP$ has an extensive derivative hierarchy.
& Generic derivative algebra plus locally uniform power-series differentiation, including
$\exp' = \exp$, $\sin'=\cos$, $\cos'=-\sin$, and hyperbolic derivatives.\\

Integration
& \gap: $\RC$ has uniform integration, properties, summation, and a fundamental theorem of
calculus.
& Radiuswise Riemann/dyadic integration, uniqueness, algebraic laws, scheme independence,
and uniform-limit preservation; no fundamental theorem located.\\

Exponential and logarithm
& \gap: $\RC$ has normalized exponentials, global exp/log laws, logarithm order and
derivative theory, rational and real powers.
& Power-series exponential with addition, positivity/order reflection/injectivity and
derivative theory.  The logarithm currently located is the local series $-\log(1-x)$ on the
open unit ball, not a global inverse to $\exp$.\\

Trigonometry
& \gap: $\RC$ has trigonometric identities, $\pi$, arcsine, circle charts/equivalences,
angle and winding infrastructure.
& Power-series sine/cosine with convergence, addition, parity, derivatives and identities;
partial tangent and local arctangent/reciprocal layers.  No $\pi$, periodicity, or global
inverse-trigonometric layer was located.\\

Hyperbolic functions
& No comparably named dedicated hierarchy was located in $\RC$ or $\RP$.
& \newcov: extensive $\sinh$, $\cosh$, $\tanh$ definitions, derivatives, addition/subtraction,
shape/order, and fundamental identities.\\

Complex analysis
& \gap: $\RC$ has complex algebra, metric completeness, modulus, inverse, polar form,
complex differentiation, Cauchy--Riemann, complex exponential/trigonometry, and roots of unity.
& No complex-number layer over $\RS$ was located.\\
\end{longtable}
\rowcolors{2}{}{}

\section{Source-to-claim ledger}

The following are the decisive checked source locations.  Paths are relative to
\path{Cubical/}.  To keep the ledger readable, $H$ abbreviates
\code{HITs/CauchyReals}, $P$ abbreviates \code{HITs/CauchyRealsP}, $G$ abbreviates
\code{Relation/Premetric/Completion}, and $S$ abbreviates
\code{Relation/Premetric/Completion/CauchySequence/Reals}.

\subsection*{Foundations and comparison}
\begin{itemize}
\item \path{H/Foundation/Base.agda}: HoTT Definition~11.3.2 constructors
  \code{rat}, \code{lim}, and \code{eqR} (Agda: $\mathsf{eq}_{\R}$).
\item \path{G/CauchySequence/Base.agda}: equations
  \eqref{eq:tail}--\eqref{eq:radiuswise} and the sequential limit constructor.
\item \path{G/Comparison.agda}: chosen, bundled, and radiuswise
  witness placements; \code{PositiveRationalChoice}; \code{SequentialTailChoice}; conditional
  comparison equivalences.
\item \path{S/Comparison/Regular.agda}:
  unconditional \code{fromRegular} and preservation of rational points/closeness.
\item \path{S/Comparison/Regular/StructuredInverse.agda}:
  \code{structuredCanonicalInverse}$\simeq$\code{sequentialTailChoice}, equation
  \eqref{eq:structured-inverse}.
\end{itemize}

\subsection*{Sequences and series}
\begin{itemize}
\item \path{H/Analysis/Sequence/Base.agda}:
  \code{IsCauchySequence'} and \code{fromCauchySequence'}.
\item \path{P/Analysis/Sequence.agda}:
  chosen-modulus \code{cauchySequenceLimit}, selected eventual order.
\item \path{S/Analysis/Sequence/Base.agda}:
  \code{chosen}$\to$\code{radiuswise}, explicit absence of the converse,
  \code{isCauchySequence}$\simeq$\code{sequenceLimit}, continuous preservation, and
  truncated eventual equality/asymptotic closeness.
\item \path{S/Analysis/Sequence/Order.agda}:
  merely eventual $\le$/$<$ and closedness under limits.
\item \path{P/Analysis/Series/Base.agda} versus
  \path{S/Analysis/Series/Base.agda}:
  one-way chosen Cauchy $\to$ summable versus
  \code{isCauchySeries}$\simeq$\code{summable}.
\item \path{P/Analysis/Series/Operations.agda} and
  \path{P/Analysis/Series/Absolute.agda} versus their $S$ counterparts:
  rational-only versus arbitrary-real scaling.
\end{itemize}

\subsection*{Majorants, products, and sharp convergence}
\begin{itemize}
\item \path{P/Analysis/Series/PowerSeries/Majorant.agda} versus the $S$
  file of the same suffix: removal of \code{NonnegativeCoefficients majorant} by
  \code{coefficientMajorantIsNonnegative}.
\item \path{P/Analysis/Series/Convolution/Mertens.agda} and
  \path{Convolution/Absolute.agda} versus
  \path{S/Analysis/Series/ConvolutionAbsolute.agda}:
  explicit sum bounds versus merely existing internal bounds; stronger absolute conclusion.
\item \path{P/Analysis/Series/Geometric/Comparison.agda} versus the $S$
  file of the same suffix: selected \code{RationalContractionBound} versus sharp $|x|<1$.
\item \path{P/Analysis/Series/PowerSeries/Convergence.agda} versus the $S$
  file of the same suffix: rational contraction clients versus the
  \code{...OfAbs<1} theorem family.
\item \path{P/Analysis/Series/Uniform.agda} versus
  \path{S/Analysis/Series/Uniform.agda}:
  selected uniform modulus versus truncated common cutoffs radius by radius.
\end{itemize}

\subsection*{Integration and transcendental extensions}
\begin{itemize}
\item \path{S/Analysis/Integration/UniformContinuity.agda}:
  $\RU(f)$ and conventional-uniform $\to$ radiuswise-uniform.
\item \path{S/Analysis/Integration/RiemannIntegral/DyadicUniformLimit/Base.agda} and its
  \path{Preservation/Global.agda}, \path{Preservation/Local.agda},
  \path{Completion/Agreement.agda}: uniform-limit continuity, preservation by the dyadic
  integral, and agreement of independent completion.
\item \code{H/Calculus/Integration/FundamentalTheorem.agda},
  \code{H/Exponentiation/}, \code{H/Trigonometry/},
  \code{H/ArcSin/}, \code{H/Circle/}, and
  \code{H/Complex/}: the principal older-coverage caveats.
\item \code{S/Analysis/Series/PowerSeries/Exponential/},
  \code{Trigonometric/}, \code{Hyperbolic/}, \code{Arctangent/}, and
  \code{Logarithm.agda}: the new $\RS$ transcendental constructions.
\end{itemize}

\section{Relation to the HoTT Book}

The local clone at \path{/Users/marcin/book} was used for terminology and notation.  The
most relevant places are:
\begin{itemize}
\item \textbf{Definition 11.2.10}: Cauchy approximations, the $\Qp$-indexed presentation used
  in the HIT constructor.
\item \textbf{\S11.3 and Definition 11.3.2}: the higher inductive-inductive Cauchy reals,
  including equations corresponding to \eqref{eq:hott-lim} and the path/extensionality
  constructor.
\item \textbf{\S11.3.4}: Cauchy completeness of the HIT real.
\item \textbf{\S3.7}: propositional truncation; the notation $\brck{A}$ used throughout this
  report follows the book.
\item \textbf{\S3.8, especially equation (3.8.1)}: the axiom of choice and the movement of
  truncation across dependent products.  Equation~\eqref{eq:choice-gap} is the relevant
  positive-rational instance.
\item \textbf{Corollary 3.9.2}: unique choice.  It does not solve the tail problem, because
  tail indices are generally not unique; only limits are unique after separatedness.
\item \textbf{\S6.9}: higher-inductive truncations, providing the conceptual background for
  eliminating merely existing estimates into propositions.
\end{itemize}

The HoTT Book's central motivation for the HIT reals is to avoid choosing representatives of
a Cauchy sequence in a quotient.  The $\RS$ development pushes the same design discipline one
level further into ordinary analysis: it avoids choosing all moduli, contraction radii, sum
bounds, uniform cutoffs, or continuity radii whenever the theorem consumes only one such
witness at a time and returns a proposition.

\section{Bottom line}

The strongest justified summary is:
\begin{enumerate}
\item The $\RS$ development accepts the radiuswise predicate
  $\prod_\varepsilon\brck{\Tail_\varepsilon}$, which is constructively weaker data than the
  old chosen-modulus predicate.
\item At exactly that strength it proves the full equivalence ``Cauchy if and only if
  convergent.''  This is the root from which most later interface strengthenings grow.
\item It systematically keeps arbitrary quantitative witnesses under truncation and opens
  them only in proposition-valued clients.  The visible payoffs are sharp $|x|<1$ hypotheses,
  bound-free Mertens, arbitrary-real scaling, uniform convergence without a global modulus,
  and radiuswise integration.
\item The missing reverse comparison is mathematically diagnosed: a structured inverse is
  equivalent to sequential tail choice, with positive-rational propositional choice as a
  sufficient assumption.
\item The new line is not yet a replacement for all of \code{HITs/CauchyReals}.  The latter
  remains materially broader in the fundamental theorem of calculus, global logarithm and
  real exponentiation, $\pi$ and inverse trigonometry, circle theory, and complex analysis.
\end{enumerate}

\begin{thebibliography}{9}
\bibitem{hottbook}
The Univalent Foundations Program,
\emph{Homotopy Type Theory: Univalent Foundations of Mathematics},
Institute for Advanced Study, 2013.
Local source consulted at \path{/Users/marcin/book}, especially Chapters 3, 6, and 11.
\end{thebibliography}

\end{document}
